\documentclass[11pt,a4paper]{article}
\usepackage[T1]{fontenc}
\usepackage{geometry}
\usepackage{setspace}
\usepackage{hyperref}
\usepackage{enumitem}

\geometry{margin=0.75in}
\setstretch{1.1}

\hypersetup{
    colorlinks=true,
    linkcolor=blue,
    pdftitle={Deborah Akuoko - Curriculum Vitae},
    pdfauthor={Deborah Akuoko}
}

\title{Deborah Akuoko - Curriculum Vitae}
\author{}
\date{}

\begin{document}

\maketitle
\thispagestyle{empty}

\section*{Profile}
Industrious self-starter looking for opportunities in academia or industry. Highly self-sufficient and experienced researcher in optical sensing, computer vision, and material detection. Able to work remotely. Excellent communication skills and ethical leadership values. Lifelong learner with genuine curiosity for advancing material-aware perception systems and autonomous systems.

Available to interview for an in-depth assessment of skills, experience and proposals.

\section*{Education}
\textbf{PhD in Engineering} - The University of Edinburgh, School of Engineering, Scotland. 2021-2025 (Expected December 2025)\\
Thesis: Material detection with SPADs\\
Advisors: Dr Istvan Gyongy (Principal), Dr Daniel Chitnis (Supporting)\\
\textit{Graduated as overall best student in class and overall best graduating Master's student, earning PhD scholarship}

\textbf{MPhil in Cancer Science} - The University of Manchester, School of Medical Sciences, England. 2019-2020\\
Thesis: Big Data for better breast cancer treatment

\textbf{MSc. Information and Communication Technology} - GIMPA School of Technology, Accra, Ghana. 2016-2017\\
Digital communication technology, Decision support systems, project management, organizational systems, ethical theories, information systems management, computer security.\\
\textit{Awarded Best Graduating Student, MSc. ICT}\\
\textit{Awarded Overall Best Student, GIMPA School of Technology, 2017}

\textbf{Bachelor of Computer Science} - The University of Cape Coast, Ghana\\
Undertaken modules in teaching and curriculum design for computer science, from graduate levels to lower primary education. Familiar with different teaching methodologies.

\section*{Work Experience}
\textbf{Machine Learning Research and Engineering} - January 2022 - Present\\
Practicing machine learning research and engineering with focus on optical sensing and computer vision. Acquired skills include:
\begin{itemize}[leftmargin=*,nosep]
    \item Remote sensing and remote depth estimation
    \item Photonic imaging technology (Single-Photon Avalanche Diode sensors)
    \item Signal processing and machine learning
    \item Deep neural networks and sequence data modeling
    \item Computer/machine vision systems
    \item Image processing and sequence data feature extraction
    \item Data processing and algorithms development
    \item Computer programming languages
    \item Machine learning frameworks (PyTorch)
    \item Project management, research data and team management
\end{itemize}
\medskip
\noindent\textit{Additional DevOps and cloud engineering skills:} Source code management (Git), Infrastructure as code (Terraform), Object-oriented programming, Deployment automation, CI/CD pipelining tools and technology (Ansible, Chef, Puppet, Jenkins), Web server technology, scripting languages (Shell, Python), Cloud infrastructure technology (Google GCP, Microsoft Azure, Amazon AWS), computer networking, systems engineering and management, Docker.

\textbf{Network and Systems Engineering with Operations} - GDS Africa, Project Assist, GHL Ghana, Accra. 2012-2018\\
Network design, implementation and troubleshooting. Experience in Layers 2 and 3 routing for CISCO and other manufacturers. Hands-on with network analytics and diagnostics. Network security and firewall policy implementation. Well-versed in network protocols: OSPF, BGP and other quality of service solutions. Skilled systems engineer with ability to adapt quickly to new systems.

\section*{Research Experience}
\textbf{PhD Research} - The University of Edinburgh, School of Engineering. 2021-2025\\
\begin{itemize}[leftmargin=*,nosep]
    \item Developed the Particle Wave Detection (PAWD) framework, formalising the dual wave-particle nature of electromagnetic energy in optical sensing systems
    \item Created novel algorithms for flat homogeneous material detection using single-pixel direct time-of-flight (dToF) histograms
    \item Developed hybrid spatiotemporal object detection system fusing SPAD transients with active pixel sensor (APS) imagery
    \item Implemented liquid purity assessment systems using time-resolved SPAD transients for non-invasive detection
    \item Built custom benchmarks and evaluation frameworks for material detection systems
    \item Developed Flask web applications for rapid evaluation and inference
    \item Created custom data formats for efficient storage of co-registered spatial and structural features
\end{itemize}

\textbf{Bioinformatics Research} - The University of Edinburgh, School of Informatics, Scotland\\
Deep learning application to antigen presentation. Data processing of genomic data.

\textbf{Clinical Breast Cancer Research} - The Christie NHS Foundation Trust. 12 months\\
Developed tumor diagnostic model for breast cancer using neural network modeling. Medical image processing, modeling patient risk of cancer treatment (radiotherapy) with tissue density, methods for estimating tissue density using Hounsfield Units, pixel manipulation and analysis.

\textbf{Bioinformatics Research} - Laboratory for Computer Science Foundations. 13 months

\section*{Research Interests}
\begin{itemize}[leftmargin=*,nosep]
    \item Signal and image processing
    \item Machine vision and computer vision
    \item Low overhead scene representation methods
    \item Safety critical computer vision
    \item Safe and reliable machine reasoning
    \item Foundation machine learning models
    \item Interactive machine learning models (e.g., LRMs)
    \item Artificial general intelligence (AGI) for Humanity
    \item Multi-modal machine learning
    \item Multi-sensor neural network modelling
    \item Spatiotemporal neural network modelling
    \item Visual object detection and tracking
    \item Learning from protected or scarce datasets
    \item Optical systems
    \item Achromatic signal processing
    \item Non-spectral signal processing
    \item Safe and beneficial artificial intelligence
    \item Material-aware perception systems
    \item Temporal signal analysis for material detection
\end{itemize}

\section*{Research Publications}
\textbf{Working Papers:}
\begin{itemize}[leftmargin=*,nosep]
    \item Remote material detection with single photon avalanche diode (SPAD) time of flight sensors (To be submitted to CVPR or other machine learning conference)
    \item Remote detection of material purity in homogenized fluid with single photon avalanche diode (SPAD) time of flight sensors (To be submitted to CVPR or other machine learning conference)
    \item Spatiotemporal neural network for improved pattern recognition (To be submitted to CVPR)
\end{itemize}

\textbf{Thesis:}
\begin{itemize}[leftmargin=*,nosep]
    \item Material detection with SPADs. PhD Thesis, University of Edinburgh, December 2025
\end{itemize}

\textbf{Previous Publications:}
\begin{itemize}[leftmargin=*,nosep]
    \item Big Data for better breast cancer treatment. MPhil Thesis, University of Manchester, 2019-2020
\end{itemize}

\section*{Awards and Scholarships}
\begin{itemize}[leftmargin=*,nosep]
    \item Winner, Rising Star of the Year, Tech and Services, Women in Tech Excellence Awards, 2020
    \item Google and Edinburgh University Scholarship, 2021
    \item Newton Fund Scholarship by DARA Big Data, 2018
    \item Best Graduating Student, MSc. ICT, GIMPA School of Technology, 2017
    \item Overall Best Student, GIMPA School of Technology, 2017
\end{itemize}

\section*{Professional Memberships and Associations}
\begin{itemize}[leftmargin=*,nosep]
    \item IEEE (Institute of Electrical and Electronic Engineering) Student Membership (since 2015)
    \item Google Professional Data Engineers (Since 2021)
\end{itemize}

\section*{Conferences and Presentations}
\begin{itemize}[leftmargin=*,nosep]
    \item International Conference on Breast Cancer Therapy, October 2021. Poster presenter, author
    \item Medical Imaging with Deep Learning (MIDL 2019)
    \item Enhancing Postgraduate Education Culture (UKCGE 2019). Session moderator
    \item IEEE PES Conference 2017, session moderator
    \item Social Innovation (SI Drive 2015). Tech Innovation session presenter
\end{itemize}

\section*{Volunteering}
\begin{itemize}[leftmargin=*,nosep]
    \item Royal Society Summer Science Exhibition, London, 2022
    \item Edinburgh University Maths Society, Maths Circle, 2022
    \item Edinburgh University Science Festival Exhibition, Edinburgh, 2022
    \item PGR Representative, Manchester University, Cancer Science Division, 2019
    \item Research Festival Representative, Manchester University, School of Medical Science, 2019
\end{itemize}

\section*{Leadership}
\begin{itemize}[leftmargin=*,nosep]
    \item Women in Technology Recruitment and Retention. Work stream lead
    \item Women in Consulting Mentorship Program. Work stream lead
    \item Girls' Prefect, Junior High School
    \item President of IAESTE Exchange Program
\end{itemize}

\section*{Teaching Experience}
\begin{itemize}[leftmargin=*,nosep]
    \item Edinburgh University Teaching Assistant (Image Processing - marking)
    \item NIIT School of Technology: Taught undergraduate students object-oriented programming concepts using Python language
    \item University of Cape Coast Model School: Taught lower primary students information technology. Designed information technology outline and adapted curriculum to suit different student learners
    \item GIMPA School of Technology Teaching Assistant: Taught graduate students in first year digital communication and signal processing systems
\end{itemize}

\section*{Teaching Interests}
\begin{itemize}[leftmargin=*,nosep]
    \item Signal processing
    \item Practical machine learning and deep neural networks
    \item Foundation machine learning models
    \item Reasoning machine learning models
    \item Practical Python programming
    \item Object-oriented programming
    \item Machine learning frameworks
    \item Multi-sensor neural network modelling
    \item Spatiotemporal neural network modelling
    \item Engineering mathematics
    \item Systems engineering
    \item Source code management and deployment automation
    \item Computer networking
    \item Data engineering pipeline development
    \item Image processing and imaging sensors
    \item Digital communication
\end{itemize}

\section*{Online Presence}
\begin{itemize}[leftmargin=*,nosep]
    \item LinkedIn: \url{https://www.linkedin.com/in/ama-d-a-minka-672ba1372/}
    \item University Profile: \url{https://eng.ed.ac.uk/about/people/ms-deborah-ewurama-akuoko}
\end{itemize}

\end{document}






